% Generated draft: requires mathematical and source review.
\documentclass[11pt]{article}
\usepackage[margin=25mm]{geometry}
\usepackage{fontspec}
\setmainfont{Latin Modern Roman}
\usepackage{amsmath,amssymb,mathtools,unicode-math}
\setmathfont{Latin Modern Math}
\usepackage{enumitem,needspace,xurl,hyperref}
\setlength{\parindent}{0pt}
\setlength{\parskip}{5pt}
\setlength{\emergencystretch}{3em}
\newcommand{\N}{\mathbb{N}}
\newcommand{\Z}{\mathbb{Z}}
\newcommand{\Q}{\mathbb{Q}}
\newcommand{\R}{\mathbb{R}}
\newcommand{\C}{\mathbb{C}}
\newcommand{\F}{\mathbb{F}}
\newcommand{\A}{\mathbb{A}}
\newcommand{\PP}{\mathbb{P}}
\newcommand{\E}{\mathbb{E}}
\newcommand{\eps}{\varepsilon}
\newcommand{\dd}{\,\mathrm{d}}
\begin{document}
\section*{UnsolvedMath ID: 30006318}
\textbf{Source problem code:} catalogue entry
\subsection{Definitions and mathematical statement}
Let \(X\) be a hyperspherical \(G\)-variety, as in the supplied source. Write \(\mathcal D\) for the formal disc and \(\mathring{\mathcal D}\) for its punctured disc, and write \(G(\mathcal D)\) and \(G(\mathring{\mathcal D})\) for the corresponding arc and loop groups. The base field, hypotheses on \(G\), and precise definition of hyperspherical are not supplied. Let \(\operatorname{Shv}(Y)\) denote the intended category of sheaves on \(Y\), and let \(\operatorname{Shv}_{G(\mathcal D)}(X(\mathring{\mathcal D}))\) denote its category of \(G(\mathcal D)\)-equivariant objects. The sheaf theory and meaning of equivariance remain unspecified.

The recoverable portion of the display attributed to Conjecture 7.5.1 is the proposed identification
\[
\operatorname{Shv}\!\left(G(\mathcal D)\backslash G(\mathring{\mathcal D})\right)
\underset{\operatorname{Shv}(G(\mathring{\mathcal D}))}{\otimes}
\operatorname{Shv}\!\left(X(\mathring{\mathcal D})\right)
\overset{?}{\simeq}
\operatorname{Shv}_{G(\mathcal D)}\!\left(X(\mathring{\mathcal D})\right).
\]
Here the tensor product is intended to be a relative tensor product of categories, using the group actions and a monoidal structure on the loop-group sheaf category; the necessary categorical setting and module structures are not specified in the record.

The supplied English statement asks for a local geometric Langlands equivalence relating this equivariant sheaf category to corresponding automorphic and spherical sheaf categories. However, the leftmost expression in the original display is corrupted, and those corresponding categories are not defined. Consequently, the displayed recoverable portion cannot be asserted to express the full intended conjecture. A precise statement of that conjecture requires recovery of the missing expression, its hypotheses, and its categorical conventions.
\subsection{Short English statement}
For a hyperspherical \(G\)-variety \(X\), the supplied question concerns a local geometric Langlands equivalence involving \(G(\mathcal D)\)-equivariant sheaves on \(X(\mathring{\mathcal D})\). A relative tensor product is recoverable from the source record, but the automorphic and spherical expression and the foundational hypotheses remain unresolved.
\subsection{Sources}
\begin{itemize}
\item Local UnsolvedMath registry, permanent ID \(30006318\), source problem code \texttt{catalogue entry}, title: \emph{Local Geometric Langlands for Hyperspherical Varieties}. The supplied record attributes the display to Conjecture 7.5.1; no individual proposer is supplied.
\item \emph{Arbeitsgemeinschaft: Relative Langlands Duality} (2025), Oberwolfach Reports, pp. 811--882, DOI 10.4171/owr/2025/18. Supplied report reference: \url{https://doi.org/10.4171/owr/2025/18}.
\item Ben-Zvi, Sakellaridis, and Venkatesh, \emph{Relative Langlands Duality}. Supplied literature reference: \url{https://www.math.ias.edu/~akshay/research/BZSVpaperV1.pdf}.
\item Literature assessment dated exactly \texttt{2026-08-22}: the supplied curation labels the problem open, identifies the surrounding source with the local geometric Langlands equivalence, and reports that no general proof was located. This is supplied dated curation, not independently verified. The background explicitly requests independent verification of time-sensitive status and cited claims.
\item The registry labels the statement \texttt{corrected\_verified} and reports that the official report was checked and notation repaired for publication. These are supplied provenance claims; this draft does not independently verify them or claim to have visited either reference.
\end{itemize}
\end{document}
